% Generated draft: requires mathematical and source review.
\documentclass[11pt]{article}
\usepackage[margin=25mm]{geometry}
\usepackage{fontspec}
\setmainfont{Latin Modern Roman}
\usepackage{amsmath,amssymb,mathtools,unicode-math}
\setmathfont{Latin Modern Math}
\usepackage{enumitem,needspace,xurl,hyperref}
\setlength{\parindent}{0pt}
\setlength{\parskip}{5pt}
\setlength{\emergencystretch}{3em}
\newcommand{\N}{\mathbb{N}}
\newcommand{\Z}{\mathbb{Z}}
\newcommand{\Q}{\mathbb{Q}}
\newcommand{\R}{\mathbb{R}}
\newcommand{\C}{\mathbb{C}}
\newcommand{\F}{\mathbb{F}}
\newcommand{\A}{\mathbb{A}}
\newcommand{\PP}{\mathbb{P}}
\newcommand{\E}{\mathbb{E}}
\newcommand{\eps}{\varepsilon}
\newcommand{\dd}{\,\mathrm{d}}
\begin{document}
\section*{UnsolvedMath ID: 30003595}
\textbf{Source problem code:} catalogue entry
\subsection{Definitions and mathematical statement}
Let \(X\) be a smooth projective variety of dimension \(n\), and let \(k\) denote its ground field. The supplied record specifies neither \(k\) nor its characteristic; no restriction on the dimension is stated. Write \(Y=X\times_k X\) and let \(\Delta_X\subset Y\) be the image of the diagonal morphism \(x\mapsto(x,x)\). Thus \(\Delta_X\) has dimension \(n\) and codimension \(n\) in \(Y\).

Under the usual numerical interpretation of bigness for cycles, let \(N_n(Y)_{\R}\) be the real vector space of algebraic cycles of dimension \(n\) modulo numerical equivalence. Define the pseudoeffective cone by
\[
\overline{\operatorname{Eff}}_n(Y)
=\overline{\left\{\sum_{i=1}^{r}a_i[Z_i]: r\geq 0,\ a_i\in\R_{\geq 0},\ Z_i\subseteq Y\text{ integral of dimension }n\right\}}.
\]
The closure and interior are taken in the ordinary topology of \(N_n(Y)_{\R}\). The diagonal is called big when
\[
[\Delta_X]\in\operatorname{int}\bigl(\overline{\operatorname{Eff}}_n(Y)\bigr).
\]
The source does not explicitly define bigness; its intended agreement with this numerical convention requires confirmation.

\textbf{Problem.} Develop geometrically useful sufficient criteria on a smooth projective variety \(X\) that imply
\[
[\Delta_X]\notin\operatorname{int}\bigl(\overline{\operatorname{Eff}}_n(X\times_k X)\bigr).
\]
This is a request for methods or sufficient hypotheses proving non-bigness in applicable cases. It does not assert that every smooth projective variety has a non-big diagonal. The source supplies no particular hypotheses for such criteria and no prescribed completeness requirement. Any field or characteristic restrictions needed by a proposed criterion must be stated explicitly.
\subsection{Short English statement}
Find useful conditions ensuring that the diagonal of a smooth projective variety is not big as a numerical cycle class on its self-product. The source asks for criteria, rather than claiming non-bigness for every variety. The ground field and the intended convention for cycle bigness require confirmation.
\subsection{Sources}
\begin{itemize}
\item Supplied local UnsolvedMath registry record, permanent ID \(30003595\), title ``Non-Bigness of Diagonals in Projective Varieties,'' source problem code \texttt{catalogue entry}, in the Oberwolfach Reports Open Problems collection. No individual proposer is supplied; the recorded proposed year is 2017.
\item Supplied report citation: \emph{Mini-Workshop: Positivity in Higher-dimensional Geometry: Higher-codimensional Cycles and Newton--Okounkov Bodies} (2018), Oberwolfach Reports, pp.~2631--2657, DOI \texttt{10.4171/owr/2017/43}. Supplied source URL: \url{https://doi.org/10.4171/owr/2017/43}. The citation's 2018 date and the record's 2017 proposed year are retained without resolving their relationship.
\item Supplied literature reference: Lehmann--Ottem, ``positivity of the diagonal,'' \url{https://sites.bc.edu/lehmann/wp-content/uploads/sites/36/2019/12/posdiag.pdf}.
\item Supplied literature reference: ``original OWR report,'' \url{https://doi.org/10.4171/OWR/2017/43}.
\item Literature checked at: \texttt{2026-08-21}. This is supplied dated curation, not an independently verified assessment. The supplied assessment labels the problem open, attributes strong necessary conditions and a classification of the nef-and-big diagonal in dimensions two and three to Lehmann--Ottem, and describes the general classification of varieties with big diagonal as subtle. It consequently regards the methodological question as unresolved in general. The scope and hypotheses of those attributed results, and the time-sensitive status, require independent verification.
\end{itemize}
\end{document}
